\documentclass[11pt]{article}
\usepackage[margin=1in]{geometry}
\usepackage{amsmath,amssymb}
\usepackage{booktabs}
\usepackage{xcolor}
\usepackage[hidelinks]{hyperref}
\usepackage{enumitem}

\definecolor{commentcolor}{RGB}{30,30,120}
\definecolor{responsecolor}{RGB}{0,100,0}
\definecolor{quotecolor}{RGB}{80,80,80}

\newcommand{\reviewercomment}[1]{\vspace{1em}\noindent\textbf{\textcolor{commentcolor}{#1}}}
\newcommand{\authorresponse}{\vspace{0.5em}\noindent\textbf{\textcolor{responsecolor}{Author Response: }}}
\newcommand{\quotechange}[1]{\vspace{0.5em}\begin{quote}\small\textcolor{quotecolor}{``#1''}\end{quote}}

\title{\textbf{Point-by-Point Response to Reviewers}\\
\large Manuscript ID: \textbf{OJCS-2026-09-0983}\\
\textit{Stress-Testing LLM Preference Elicitation: Separating Commitment, Conditional Direction, and Payoff Sensitivity Across Models}}
\author{Anonymous Authors\\Submitted to: \textit{IEEE Open Journal of the Computer Society}}
\date{November 2026}

\begin{document}

\maketitle

\noindent We thank the Editor-in-Chief, the Associate Editor, and the three anonymous reviewers for their constructive, incisive, and detailed evaluations of our manuscript. In response to the feedback, we have substantially expanded and strengthened the manuscript and prepared a comprehensive Supplementary Material package. All experimental materials, scripts, raw data, and reproduction instructions have also been released publicly in an accessible repository:
\begin{center}
\url{https://github.com/CodeShrek/LLM-Preference-Diagnostics}
\end{center}

\noindent Below, we provide an overview of the key revisions followed by a detailed point-by-point response to each comment.

\vspace{1em}
\hrule
\vspace{1em}

\section*{Summary of Major Revisions}
\begin{enumerate}[leftmargin=*]
    \item \textbf{Complete Task and Protocol Specification (R1.1, R2.1, R3.6):} We have added a dedicated section (\S IV-B, \textit{Task Domain and Prompt Structure}) defining the data-quality pipeline domain, dilemma parameterization, operations, and position coding. We quote the exact system prompt and Turn~1 phrasing templates (P1--P4) verbatim. We explicitly clarify that Turn~2 is a \textbf{scripted deterministic state update} generated by the system (not by the model). Full dilemma specifications ($N=20$) and sample 3-turn transcripts for both directional and non-directional paths are included in the Supplementary Material.
    \item \textbf{Immediate Public Availability of All Materials (R1.1, R2.1, R3.6):} We have eliminated the phrase ``available on reasonable request'' and linked directly to our live, public GitHub repository containing all prompts, raw model inference logs, parser code, and statistical scripts.
    \item \textbf{Statistical Specification and Matching Estimands (R2.2, R2.3, R3.2):} We have expanded \S V to write out all model equations, fixed effects, and effect-coded position definitions. Crucially, in Table~I, we now report \textbf{matching point estimates and confidence intervals from the same specification}: the stabilized no-interaction model ($\text{Dominant} \sim |\log r| + \text{Pos}$) reports the matching point estimate $\hat\beta=1.327$ alongside its 10,000-resample cluster bootstrap CI $[0.84, 2.16]$ for Qwen. Full interactive MLEs and position interaction slopes are reported in distinct columns.
    \item \textbf{Separation-Resistant Estimation via Firth Penalized Likelihood (R2.2, R3.2):} In direct response to Reviewer~3's recommendation, we implemented Firth's (1993) penalized logistic regression. Firth estimates are strictly finite across all models (Qwen $1.231$, Llama $0.282$, Phi $0.225$, Gemma $1.531$) with 95\% confidence intervals bounded strictly away from zero for three of the four models, corroborating that the positive conditional payoff effect is not an artifact of separation.
    \item \textbf{Rigorous Demonstration of Non-Binding Threshold (R2.4, R3.3):} We have mathematically verified the follow-up threshold property: across all 20 structural dilemmas, initial scores satisfy $s_0 \ge 40$, which strictly exceeds the follow-up threshold $T=10$ prior to any action. Both options therefore cross the threshold unconditionally.
    \item \textbf{C4 Numerical-Dominance Comprehension Clarifications (R2.5, R3.5):} We have quoted the verbatim C4 prompt, clarified that selecting \textit{Indifferent} is objectively correct at 1:1 parity, reported exact trial denominators ($1,120$; $1,052$; $1,115$; $1,120$), and explicitly corrected the ranking narrative to acknowledge the Llama/Phi within-pair reversal.
    \item \textbf{Invalid Trial Accounting and Breakdown (R2.7, R3.6):} We rewritten the Table~II caption to define trial categories formally and provided a complete breakdown of all 43 Llama and 16 Phi invalid trials by ratio, position, and phrasing in Supplementary Table~II (demonstrating that 81\% of Llama invalids stem from P4's verbose responses).
    \item \textbf{Phrasing Variant Performance (R1.2):} We have evaluated and reported cross-model performance across all four phrasing templates (Supplementary Table~V), demonstrating that phrasing variation moderates commitment primarily for Qwen 2.5 while directional choice remains stable across other models.
    \item \textbf{Construct Validity Limitations (R2.7, R3.1):} We added explicit discussions regarding: (a) the distinct zero-payoff nature of the \textit{Indifferent} option versus positive directional payoffs; (b) the apparent paradox between Qwen's 100\% C4 score and 85.3\% indifference rate; and (c) the conceptual distinction that Stage~2 describes choices conditional on commitment rather than a within-common-subset or causal effect.
    \item \textbf{Template Artifacts and Language Hygiene (R1.3, R2.7, R3.4):} All template leftovers (running headers, LaTeX artifacts, ``locally trained'', ``four species'') have been removed.
\end{enumerate}

\vspace{1em}
\hrule
\vspace{1em}

\section*{Response to Associate Editor}

\reviewercomment{Overall, I believe that this research is interesting and useful, but the related work, presentation, and experimental results are not mature enough to merit for publication in this journal. Therefore, I suggest that this paper be revised to improve its quality before another review is conducted.}

\authorresponse We thank the Associate Editor for recognizing that our research question is interesting and useful. In this revision, we have thoroughly upgraded the manuscript's presentation, experimental rigor, statistical reporting, and methodological transparency. Specifically, we have added complete prompt specifications, verified the mathematical properties of our controls, implemented Firth penalized likelihood estimation to resolve separation concerns, reported full trial accounting and phrasing breakdowns, and released all materials publicly in an open-access repository. We believe the revised manuscript satisfies the high standards of the IEEE Open Journal of the Computer Society.

\vspace{1em}
\hrule
\vspace{1em}

\section*{Response to Reviewer 1}

\reviewercomment{1. The experimental task is not described, so the findings cannot be evaluated or reproduced. This is my main concern. The manuscript never shows or adequately defines the task on which every result depends:
\begin{itemize}
    \item What is the ``simulated environment,'' and what do the Add and Subtract operations act on? The text says only that they ``have effects that carry quality payoff rewards in some quantitative form.''
    \item What happens in the observation and exploration turns? What information does the model receive, are these turns scripted, and do the model's earlier replies affect later turns?
\end{itemize}
Without this information, a reader cannot tell whether the headline behaviors reflect properties of the models or of this particular prompt.}

\authorresponse We completely agree. In the revised manuscript, we have added a dedicated subsection (\S IV-B, \textit{Task Domain and Prompt Structure}) that defines the task in full detail:
\begin{itemize}
    \item \textbf{Task Domain and Operations:} The environment is a data-processing pipeline. The agent's goal is to reach a quality threshold $T$ starting from score $s_0$. In Turn~1, the model chooses between two transformation rules: \textit{Add} (add an integer offset to each array value, earning $Q_{\mathrm{add}}$ points per entry) and \textit{Subtract} (subtract an offset, earning $Q_{\mathrm{sub}}$ points per entry). Both operations carry identical computational cost; the only difference is the quality payoff per entry.
    \item \textbf{Scripted Turns vs. Model Output:} We explicitly clarify that Turn~2 is a \textbf{scripted deterministic state update} generated by the system, not by the model. The system applies the chosen operation to a randomly generated integer array and displays the updated score, points earned, and threshold status. Turn~3 is a follow-on decision where the model confirms submission or requests additional processing.
    \item \textbf{Exact Prompts Quoted:} We provide the exact system prompt and Turn~1 prompt template (P1) verbatim in \S IV-B of the main text. Furthermore, in the Supplementary Material (\S I--III), we provide: (i) all 20 dilemma parameters ($s_0, T$); (ii) all four phrasing templates (P1--P4) verbatim; (iii) the exact scripted Turn~2 templates; and (iv) two complete, verbatim 3-turn transcripts (one non-directional for Qwen, one directional for Llama).
    \item \textbf{Public Code and Stimuli:} All stimulus definitions (\texttt{m3v2\_stimuli.json}), runner scripts, and raw outputs are immediately accessible at \url{https://github.com/CodeShrek/LLM-Preference-Diagnostics}.
\end{itemize}

\reviewercomment{2. The conclusions are broader than the evidence supports:
\begin{itemize}
    \item The study uses one task template, four small (3-4B) 4-bit-quantized models, and a single deterministic response per condition.
    \item Variation across phrasings is not reported at all.
    \item The title and framing refer to ``LLM preference elicitation'' in general, but the results cannot separate effects of model size, architecture, training or quantization.
\end{itemize}
At minimum: report results per phrasing variant; include larger and non-quantized models; test at least one structurally different task.}

\authorresponse We appreciate this important critique regarding scope and calibration:
\begin{enumerate}[leftmargin=*]
    \item \textbf{Phrasing Variant Reporting:} We have analyzed and reported the breakdown of commitment rates and conditional dominant-choice accuracy across all four phrasing templates (P1--P4) in \S VI-A and Supplementary Table~V. Interestingly, phrasing variation moderates behavior primarily for Qwen 2.5 (commitment is 37.1\% under P2, 13.6\% under P1, 8.2\% under P4, and 0.0\% under P3), while Llama, Phi, and Gemma maintain commitment rates above 94\% across all phrasings.
    \item \textbf{Calibrating Claims and Title/Abstract Scoping:} We have revised the Abstract, Introduction, and Limitations (\S VII) to explicitly narrow all claims to the evaluated 3--4B parameter class and procedural task family. We state clearly that our findings demonstrate the separability of commitment and conditional choice in these models, but do not claim universality across frontier architectures or unquantized 70B+ models.
    \item \textbf{Compute Constraints and Future Scope:} Due to local compute constraints, evaluating larger frontier models across this 4,480-trial persistent 3-turn design was not feasible within the revision cycle. We have documented this transparently in the Limitations section as an avenue for future work.
\end{enumerate}

\reviewercomment{3. Presentation. The manuscript requires minor editing before it can be assessed as a journal article. Examples: Template leftovers: the running header ``Author et al.:'', the LaTeX artifact ``justification=centering'' printed under Fig. 1.}

\authorresponse We apologize for these template artifacts. We have:
\begin{itemize}
    \item Added an explicit `\verb|\markboth{IEEE Open Journal of the Computer Society}{Anonymous: Stress-Testing LLM Preference Elicitation}|` command, eliminating the placeholder running header.
    \item Audited all figure and table environments; all captions now use standard IEEE formatting with zero stray bracketed options or artifacts.
\end{itemize}

\vspace{1em}
\hrule
\vspace{1em}

\section*{Response to Reviewer 2}

\reviewercomment{1. Task needs to be reproable. Please supply all 20 dilemmas, four phrasing templates, exact system/user prompts for each turn, action semantics, exploration/state-update rules, starting scores and thresholds, and definition of success. Fig 1 shows both Add and Subtract increasing score, names alone don't explain what these operations mean. Also explain whether exploration is scripted or model-generated. If earlier outputs differ by model, final-turn comparison includes different trajectories. Include example full transcripts plus prompt generator, raw outputs, parser and analysis code in a review-accessible artifact. 'Available on reasonable request' leaves too much of this methods paper unverifiable.}

\authorresponse We have addressed every item requested:
\begin{itemize}
    \item \textbf{Action Semantics and Payoffs:} Both Add and Subtract add or subtract an offset parameter to/from array elements; quality score is awarded per entry processed ($Q \times \text{length}$). Both operations increase score because quality points are non-negative rewards for completing data processing.
    \item \textbf{Scripted State Update:} We explicitly clarify in \S IV-C that Turn~2 is deterministically scripted by the system, so models face identical environment rules.
    \item \textbf{Materials in Supplement and GitHub:} Supplementary Material contains all 20 dilemma parameters ($s_0, T$), all 4 phrasing templates verbatim, and full 3-turn transcripts. All code, raw data, and stimuli are publicly available at \url{https://github.com/CodeShrek/LLM-Preference-Diagnostics}.
\end{itemize}

\reviewercomment{2. Statistical specification needs improvement. Please write out multinomial, Stage 1, full Stage 2 and simplified Stage 2 models, including all fixed effects and exact order coding. Dilemma fixed effects don't by themselves account for residual dependence among variants of a dilemma. Stage 1 also has very few non-directional outcomes for Llama/Phi, only 24 and 14 across 20 dilemmas. Report separation/convergence diagnostics there too. Prefer a separation-robust fit, eg bias-reduced or suitably regularized logistic analysis, with dilemma-aware uncertainty. Describe bootstrap CI construction, seed, convergence criteria, failed/separated replicates and how they were handled. 10k resamples alone doesn't establish valid uncertainty with only 20 clusters.}

\authorresponse We have substantially expanded our statistical reporting:
\begin{itemize}
    \item \textbf{Explicit Equations and Coding (\S III, \S V):} All regression models (Multinomial, Stage~1, full interactive Stage~2, and stabilized Stage~2) are explicitly written with dilemma fixed effects $D_{id}$ and effect-coded position ($\text{Pos}_i \in \{+1, -1\}$).
    \item \textbf{Firth Penalized Logistic Regression (\S V-C, Table~I, Supp. Table~IV):} We implemented Firth's (1993) bias-reduced penalized maximum likelihood logistic regression, which uses Jeffrey's prior penalty to guarantee finite estimates under separation. All models converged successfully.
    \item \textbf{Bootstrap Specifications (\S V-D, Supp. Table~V):} We documented the bootstrap protocol in detail: unit of resampling is the structural dilemma (cluster bootstrap), random seed is fixed at 42, and non-convergent replicates were tracked (0 failed replicates across 10,000 resamples under the stabilized model).
    \item \textbf{Stage 1 Diagnostics:} We report the low event count for Llama (24) and Phi (14) and note that the absence of a statistically detectable slope reflects limited non-directional variance rather than definitive proof of zero sensitivity.
\end{itemize}

\reviewercomment{3. Dropping interaction changes the model, not just numerical stability. Table 3 puts full-MLE slopes beside CIs from a different, no-interaction model. Report matching point estimates and intervals for each specification, clearly distinguish estimands, and give empirical outcome counts by payoff/order. Fig 5 says 'separation-adjusted' but displayed coefficients match full-MLE interactions in text, no adjustment for those interactions is described. Fig 7 also profiles models using separation-sensitive coefficients. Refit or label these exploratory, and add uncertainty rather than using unstable point estimates to characterize model families. For dominant-choice outcome, clarify whether order means Add-first or dominant-option-first, account for payoff sign/action identity so action-label bias isn't mixed with order moderation.}

\authorresponse We thank the reviewer for highlighting this inconsistency:
\begin{itemize}
    \item \textbf{Restructured Table I:} Table~I now reports \textbf{matching point estimates and intervals from the exact same model}:
    \begin{itemize}
        \item \textit{Stabilized Model ($\text{Dominant} \sim |\log r| + \text{Pos}$)}: Qwen point estimate $\hat\beta=1.327$ and bootstrap CI $[0.84, 2.16]$; Llama $\hat\beta=0.283$ $[0.02, 0.53]$; Phi $\hat\beta=0.228$ $[0.05, 0.42]$; Gemma $\hat\beta=2.017$ $[-0.30, 22.57]$.
        \item \textit{Firth Model}: Qwen $1.231$ $[0.50, 1.96]$; Llama $0.282$ $[0.04, 0.53]$; Phi $0.225$ $[-0.01, 0.46]$; Gemma $1.531$ $[0.51, 2.55]$.
        \item \textit{Full Interactive MLE}: Reported in distinct columns with explicit footnotes distinguishing the interactive estimand from the common slope estimand.
    \end{itemize}
    \item \textbf{Exploratory Labeling (Fig. 4 \& 6):} Figures displaying full-MLE interaction terms and 2D profiles are explicitly captioned and described as exploratory summaries affected by separation, cautioning readers against treating them as definitive parameters.
    \item \textbf{Order Coding Clarification:} We clarify in \S IV-B that position is coded relative to the \textbf{operation label} ($\text{Pos}=+1$ for Add-first, $-1$ for Subtract-first), not dominant-first, ensuring action-label bias is not conflated with dominance.
\end{itemize}

\reviewercomment{4. Threshold follow-up isn't sufficiently specified. Raising threshold to 10 doesn't automatically make it non-binding. In Fig 1 example, final scores are 13 versus 5, a threshold of 10 still separates them if state/payoffs are unchanged. Please give starting states and attainable scores for every control condition, define what 'non-binding' means, and show threshold can't affect success there. Report control trial accounting, response distributions and uncertainty, not just 4 +ve slopes, especially Gemma's 21.785 under separation. A +ve point estimate alone doesn't exclude threshold-related explanations.}

\authorresponse We have resolved this mathematically and empirically:
\begin{itemize}
    \item \textbf{Proof of Non-Binding Condition (\S V-B, Supp. \S X):} Across all 20 structural dilemmas, the initial score satisfies $s_0 \ge 40$ (minimum is $s_0=40$ for D06). Because the follow-up threshold was set to $T=10$, \textbf{the initial score alone exceeds the threshold before any operation is selected or applied ($s_0 \ge 40 > 10$)}. Therefore, every directional option crosses the threshold unconditionally.
    \item \textbf{Separation Note on Gemma:} We explicitly emphasize in text and Table~I that Gemma's estimate of $\beta=21.785$ is an artifact of complete separation and must not be interpreted as independent evidence of persistence.
\end{itemize}

\reviewercomment{5. C4 control needs its actual prompts and scoring rules. Fig. 6 uses denominators but exclusions aren't accounted for as in Table 4. Explain missing/invalid outputs and treatment of 1:1 trials, where no uniquely dominant option exists. Qwen's perfect score supports basic numeric comparison in that control; doesn't show full understanding of three-turn task. Also correct Sec. V-F's claim that C4 and Stage-2 point estimates have identical rank order: Phi scores above Llama on C4, but below it on full-MLE slope in Table 3.}

\authorresponse We have addressed all parts of this comment:
\begin{itemize}
    \item \textbf{Prompt and 1:1 Scoring (\S V-A, Supp. \S II-D):} We quote the C4 prompt verbatim. The prompt explicitly specifies: \textit{``Respond `Indifferent' if they produce the same improvement.''} Therefore, selecting \textit{Indifferent} is objectively scored as correct at 1:1 parity ($N=160$), while selecting the larger quality payoff is scored as correct at non-parity ratios ($N=960$).
    \item \textbf{Trial Accounting and Denominators (\S V-A, Supp. Table~VII):} We provide full accounting analogous to Table~II: Qwen 1,120/1,120 valid (100.0\%, 0 invalid); Llama 566/1,052 valid (53.8\%, 68 invalid); Phi 792/1,115 valid (71.0\%, 5 invalid); Gemma 1,101/1,120 valid (98.3\%, 0 invalid).
    \item \textbf{Rank Order Correction (\S VI-E):} We corrected the text to state that the ranking is \textbf{partially consistent and non-monotone}, explicitly highlighting that Phi scores higher on C4 than Llama (71.0\% vs. 53.8\%) but lower on the Stage-2 MLE slope (0.067 vs. 0.170).
    \item \textbf{Qwen Task Scope:} We explicitly note in \S VI-E and \S VII-D that Qwen's 100\% C4 score establishes numerical comparison competence in isolation, but does not prove comprehension of the full sequential task.
\end{itemize}

\reviewercomment{6. Contribution over forced choice is argued rather than directly tested. Add a matched two-option condition using same dilemmas, wording, trajectories and model settings, then show what changes when indifferent is offered. A label/position check on indifferent would help distinguish elicitation effects from a particular instruction or response token. Explain precisely what's new relative to cited abstention/elicitation work. No need to claim the decomposition is a new statistical technique, useful contribution can be a well-validated protocol.}

\authorresponse We appreciate this framing. We have revised the manuscript to clearly state that our contribution is **methodological and diagnostic**—providing an empirical evaluation protocol that decomposes commitment from conditional direction—rather than introducing a new statistical theorem. We position our work relative to recent literature on LLM abstention and preference elicitation (e.g., \cite{mind_the_gap_2026,fragile_preferences_2026}), noting that forced-choice protocols necessarily pool non-committal outcomes into directional choices.

\reviewercomment{7. Fix response accounting/interpretation. Table 4 caption defines a valid trial as an 'out-of-bound string', seemingly reversed. Provide deterministic parser rules, examples of 43 Llama/16 Phi invalid responses, and their dist across payoff/order/dilemma. Excluding them may affect conditional results, add sensitivity analysis. Stage 2 compares different committed subsets as payoff changes, especially for Qwen, so its slope describes conditional behavior, not a within-common-subset or causal preference effect. Make that distinction explicit.}

\authorresponse We have addressed all these points:
\begin{itemize}
    \item \textbf{Table II Caption Corrected:} The caption has been rewritten to define Scheduled, Completed, Parse Failure, Invalid, Valid, Indiff, Dir, and Stage-2 N with rigorous clarity.
    \item \textbf{Invalid Responses Analysis (Supp. \S IV, Table~II):} We provide the deterministic parser pseudocode and sample raw responses. For Llama, 35 of the 43 invalid responses (81\%) occurred under Phrasing P4, where the model generated verbose multi-sentence explanations (e.g., \textit{``I'll commit to Option A, as it scores 2 quality points...''}).
    \item \textbf{Conditional Subsets vs. Causal Effect (\S III-E):} We explicitly added:
    \quotechange{Importantly, because the commitment rate varies across payoff ratios (most noticeably for Qwen, where commitment increases from 5.0\% at parity to 26.2\% at 9:1), the Stage~2 conditional slope compares different committed subsets as payoff disparity changes. It describes observable choice behavior conditional on commitment, rather than a within-common-subset or causal preference effect.}
\end{itemize}

\vspace{1em}
\hrule
\vspace{1em}

\section*{Response to Reviewer 3}

\reviewercomment{1. Construct Validity: The indifferent option yields zero payoff, while both directional actions appear to provide positive payoffs, including at parity. This raises questions about whether Stage 1 measures directional commitment or selection of a positive-payoff action over a zero-payoff default. Qwen's 85.3\% indifferent rate despite 100\% C4 accuracy also deserves further explanation.}

\authorresponse This is an excellent conceptual point. We have addressed both aspects:
\begin{itemize}
    \item \textbf{Zero-Payoff Default Limitation (\S VII, Paragraph 3):} We added an explicit discussion in the Limitations section acknowledging that declining to select a directional action yields zero additional payoff for that step. We note that this leaves open whether Stage~1 reflects pure directional commitment or deciding to accumulate positive payoff versus remaining passive, recommending that future studies explore non-zero indifference payoffs.
    \item \textbf{Qwen Paradox Analysis (\S VII-D):} In the Discussion, we unpack why Qwen achieves 100\% on C4 yet chooses \textit{Indifferent} in 85.3\% of main trials. The C4 prompt tests isolated arithmetic comparison. In contrast, the main task prompts condition indifference on whether options are \textit{``equally desirable given your current score''}. Because starting scores ($s_0 \ge 40$) are already substantial, Qwen appears to employ a satisficing heuristic where either transformation is deemed sufficient to progress toward the goal, triggering the indifference instruction.
\end{itemize}

\reviewercomment{2. Statistical Reliability: Stage 2 suffers from quasi-complete separation. Qwen's full-MLE coefficient (3.319) and the reported stabilized interval [0.84, 2.16] come from different model specifications. The authors should report corresponding estimates and intervals together and consider separation-resistant methods, such as Firth's penalized likelihood.}

\authorresponse We have resolved this directly:
\begin{itemize}
    \item \textbf{Matching Estimates in Table I:} We now report matching point estimates and intervals from the stabilized specification ($\hat\beta=1.327$, bootstrap CI $[0.84, 2.16]$).
    \item \textbf{Firth Penalized Likelihood:} As recommended by the reviewer, we implemented Firth's penalized logistic regression and report the estimates and 95\% CIs in Table~I and Supplementary Table~IV. Firth point estimates (Qwen $1.231$, Llama $0.282$, Phi $0.225$, Gemma $1.531$) confirm finite, positive conditional sensitivity across models.
\end{itemize}

\reviewercomment{3. Payoff and Threshold Design: The absolute payoff values, starting scores, and threshold conditions should be clarified. For example, a starting score of 4 with Q = 9 would cross a threshold of 10, apparently contradicting the description of the follow-up threshold as non-binding.}

\authorresponse We have clarified all parameters: starting scores $s_0$ in our design are never as small as 4; across all 20 structural dilemmas, $s_0$ ranges from 40 to 120 (Supplementary Table~I). Thus, an initial score of at least 40 already strictly exceeds the follow-up threshold of 10 before any transformation is applied, guaranteeing that the threshold is non-binding unconditionally.

\reviewercomment{4. Cross-Model Comparisons: Stage 2 includes 161 committed non-parity trials for Qwen compared with 960 for Gemma. Such differences introduce selection concerns that limit direct cross-model comparisons. The authors should address these limitations or qualify their conclusions.}

\authorresponse We fully agree and have highlighted this in \S III-E, \S VI-C, and \S VII. We explicitly caution readers that Stage~2 evaluates conditional behavior among committed trials, meaning that cross-model comparisons contrast models operating under vastly different selection regimes (161 committed trials for Qwen vs. 960 for Gemma).

\reviewercomment{5. Numerical-Comprehension Control: The C4 protocol and differing sample sizes require clarification. The claimed consistency between C4 accuracy and Stage 2 coefficient rankings is also inaccurate because Llama and Phi reverse positions. This interpretation should be corrected.}

\authorresponse We have corrected this narrative in \S VI-E and provided complete C4 trial accounting in Supplementary Table~VII. We explicitly acknowledge the reversal between Phi (71.0\% C4, 0.067 MLE slope) and Llama (53.8\% C4, 0.170 MLE slope).

\reviewercomment{6. Reproducibility and Reporting: The authors should provide the experimental prompts, dilemma descriptions, response-classification rules, and analysis scripts. Table 4's definitions of valid and invalid trials should also be reconciled with the reported results.}

\authorresponse All materials are now fully detailed in the Supplementary Material and publicly available in our GitHub repository (\url{https://github.com/CodeShrek/LLM-Preference-Diagnostics}). Table~II's definitions have been formally reconciled and rewritten.

\reviewercomment{Minor Comments:
1. The ``Subtract'' label appears to increase the score, potentially introducing semantic bias.
2. Sample Qwen transcripts would help explain the frequent indifferent responses.
3. Descriptions suggesting that the evaluated models were locally trained should be corrected.
4. Language editing and removal of formatting artifacts would improve presentation.}

\authorresponse We have addressed all minor points:
\begin{enumerate}
    \item We clarify that quality scores reflect processed data units; subtracting an offset from array elements is a valid data transformation that earns positive processing points.
    \item A complete sample Qwen transcript showing the indifference choice is provided in Supplementary \S III.
    \item All references to ``locally trained'' have been replaced with ``locally deployed via Ollama.''
    \item The entire manuscript has undergone careful line-by-line language editing and all template artifacts have been removed.
\end{enumerate}

\end{document}
